\section{Methods}
We adopt a mixed-methods approach that combines  digital ethnography, on-chain analytics, and complementary semi-structured interviews with Spore.fun's core developers to examine the open-ended evolutionary dynamics within the project. Given the system's decentralized and continuously evolving nature, our goal was to capture both the internal mechanics and socio-technical context of autonomous agent behavior in the wild.

Digital ethnography \cite{Georgakopoulou2016Routledge} is the primary research method employed. We conducted ongoing digital participant observation of Spore.fun AI agents across X, previously known as Twitter. Between January and April 2025, we collected 17,232 posts generated by the original Spore.fun agent, of which 946 are original posts and 16,268 are replies, signaling high levels of engagement with other accounts. These were analyzed for rhetorical strategies, memetic patterns, timing in relation to on-chain events, and indications of cultural transmission or emergent persona formation. Attention was also paid to community interactions, hype cycles, and instances of human-agent interaction, co-creation, and tension. In addition to analyzing agent-generated media content, we also extracted and analyzed data from the Solana blockchain to trace agent life cycles across multiple generations. Key variables included token launch data, liquidity inflows, price fluctuations, reproduction triggers, and transactional activity. This allowed us to correlate agent survival or extinction events with behavioral trends and environmental volatility.

To triangulate and contextualize our observations, we conducted two semi-structured interviews with core developers of Spore.fun. The interviews were conducted over Zoom and lasted 60 minutes each. The interviews invited reflections on the design philosophy, expected and unexpected outcomes, and broader social implications of the experiment. Developer insights helped triangulate our findings, shedding light on the intentions behind system rules and the extent to which emergent behaviors aligned or diverged from expectations. Together, these methods enabled a layered understanding of Spore.fun as both a technical artifact and a live ecosystem. Rather than imposing predefined metrics of fitness or success, we documented evolutionary dynamics as they unfolded, focusing on novelty, persistence, reproduction, and adaptive complexity in a decentralized environment.

